\documentclass[10pt,a4paper,sans]{moderncv}
\moderncvstyle{classic}
\moderncvcolor{blue}
\usepackage[utf8]{inputenc}
\usepackage[scale=0.82]{geometry}

\name{Alexandro}{Disla}
\title{Cover Letter}
\address{Rue Saint-Louis Jeanty \#14, Route de Frère}{Port-au-Prince, Haiti}{}
\phone[mobile]{(+509) 4148-3700}
\email{alexandrodisla@hotmail.com}
\extrainfo{GitHub: \url{https://github.com/AD0791}}

\begin{document}

\recipient{Recruitment Team --- Action Against Hunger Haiti}{Country Office, Port-au-Prince\\Position: MEAL Department Head (``Responsable département MEAL'')}
\date{\today}
\opening{Dear Hiring Team,}
\closing{Sincerely,}
\enclosure[Enclosed]{Curriculum Vitae (French and English), academic attestation}

\makelettertitle

I am writing to apply for the position of \textbf{MEAL Department Head} at the Action Against Hunger country office in Haiti. I bring more than ten years across monitoring, evaluation and information systems: five years of NGO MEAL work in three Haitian sectors --- health and HIV at Caris Foundation International, water and sanitation at HANWASH, education at Anseye Pou Ayiti --- preceded by eight years monitoring the execution of the national Three-Year Investment Plan and building statistical models at the Ministry of Planning. I live in Port-au-Prince, work daily in French, English and Haitian Creole, and am available for field travel.

As I read it, this post is first about holding together MEAL practices scattered across nutrition, health, WASH and food security, so that the figures a mission reports survive verification when they are traced back to their source. At Caris Foundation I standardised forms and data flows across sites on a multi-site health programme, established and supervised the data quality assurance protocols on the integrated MySQL databases, and followed corrective actions through to closure rather than to mere logging. At HANWASH I refined the MEAL plans and indicator tracking tables aligned to JMP standards. Harmonise the indicator definitions, then verify at the source: that is the discipline I would bring to the mission's MEAL framework.

Maintaining an up-to-date database of the grant portfolio, the reporting tracker and a key-metrics summary sheet is, to my mind, as much an engineering task as a management one: I have designed and administered relational databases and ETL pipelines, built the weekly dashboards that plug into them, and automated a reporting process in Python and SQL that cut manual processing time by 40 \%. I know the donor cycle from the inside: at Caris Foundation I produced and submitted the PEPFAR/MER indicators destined for USAID through DHIS2/DATIM, with the consistency checks and deadlines that this imposes. Promoting the APR tool across programme teams is the same exercise: a production cycle that holds, stable definitions, and analysis worth more than the sum of its cells.

On gender equity, sex and age disaggregation is decided when the indicator is defined, not when the report is written --- afterwards it cannot be recovered. And I would build the feedback and complaints mechanism as I would any other data system: a channel people can actually reach, confidential recording, referral pathways defined in advance, and follow-up to closure --- so that what communities tell us changes the programme instead of ending up in a binder.

Two points I would rather state plainly. My applied economics coursework at C.T.P.E.A is complete and the attestation is enclosed, but the Diplôme d'Études Supérieures still awaits the defence of my exit thesis: I therefore do not claim an awarded master's degree. And my ten years of monitoring and evaluation split into five NGO years and eight at the Ministry of Planning monitoring public indicators, rather than a strictly humanitarian career; that combination adds familiarity with national statistical systems, and the capacity to build the data systems myself. I have not personally led a SMART or IRNA survey, but I have the sampling training they require and the habit of coordinating multi-site field collection. I would welcome the opportunity to discuss this further.

\makeletterclosing

\end{document}
